\section{预备知识与问题定义}

\subsection{RTL 安全属性与转移系统}

给定 RTL 设计 $D$、目标安全属性及环境假设，形式化后端依据统一的时钟、复位和 SVA 语义构造状态转移系统 $T=(V,\mathrm{Init},\mathrm{Tr})$。其中，$V$ 表示当前状态变量，$V'$ 表示其后继状态变量；$\mathrm{Init}(V)$ 刻画合法初始状态，$\mathrm{Tr}(V,V')$ 刻画合法的一步转移。本文用状态谓词 $p(V)$ 表示展开后的目标安全条件，并把从初始状态经有限次合法转移可到达的状态集合记为 $\mathrm{Reach}(T)$。证明目标是 $T\models AGp$，即对每个 $s\in\mathrm{Reach}(T)$ 均有 $p(s)$ 成立。对含时序运算的属性，上述状态记号指后端按相同语义展开后的模型，而不是直接对原始 SVA 文本作布尔求值~\cite{peled2026large}。

有界模型检查（BMC）从合法初始状态搜索长度不超过给定深度 $d$ 的执行轨迹，可用于发现违反 $p$ 的可达反例。未找到反例只能说明该深度内未发现目标违例，不能据此报告 \texttt{PROVED}；证明成功仍须依靠下文所述的初始性与归纳保持检查。

\subsection{归纳增强与形式化检查}

状态谓词 $I(V)$ 是归纳不变式，当且仅当它满足初始性和归纳保持：
\[
\begin{array}{l}
\mathrm{Init}(V)\Rightarrow I(V),\\
I(V)\land\mathrm{Tr}(V,V')\Rightarrow I(V').
\end{array}
\]
若还满足 $I(V)\Rightarrow p(V)$，它便能证明目标安全属性。目标 $p$ 本身可能不满足归纳保持；此时可用辅助引理组成归纳增强。对任意候选集合 $S$，定义
\[
I_S(V)=p(V)\land\bigwedge_{\ell\in S}\ell(V),
\]
其中空集合上的合取取真，故 $I_{\emptyset}=p$。所选引理无需各自构成归纳不变式；只有整个 $I_S$ 满足上述两项条件，$S$ 才是目标 $p$ 的有效归纳增强。

实际检查可表述为两项可满足性查询：
\[
\begin{array}{rl}
Q_{\mathrm{init}}(S)={}&\mathrm{Init}(V)\land\neg I_S(V),\\
Q_{\mathrm{step}}(S)={}&I_S(V)\land\mathrm{Tr}(V,V')\\
&{}\land\neg I_S(V').
\end{array}
\]
两项均为 UNSAT 才能返回证明；$Q_{\mathrm{step}}$ 的 SAT 见证可能始于不可达状态，不能直接视为目标属性的可达反例。$Q_{\mathrm{init}}$ 为 SAT 时也须区分是 $p$ 还是某条候选在初始状态失败。TIMEOUT、UNKNOWN 或工具错误均不构成证明。对无法按后端语义可靠判定单步反例中候选真值的时序 SVA，反馈求值记为 unknown，不据此作确定性剪枝。

\subsection{固定候选池上的引理选择问题}

候选生成结束后，固定引理池 $C=\{\ell_1,\ldots,\ell_n\}$，并为每条候选分配不变的 ID；候选内容在选择阶段保持不变。给定同一设计 $D$、目标 $p$、环境假设、后端语义及候选池 $C$，选择器安排待检验子集 $S\subseteq C$ 的顺序；它不生成新引理，也不修改设计或目标属性。候选仅表示待验证的性质，并不预设其满足初始性、归纳保持或与目标相关。下文公式中的 $S$ 表示引理集合，实际输出则由其 ID 列表唯一表示。

固定顺序的前缀搜索只访问以下组合：
\[
P_k=\{\ell_1,\ldots,\ell_k\},\qquad 1\le k\le n.
\]
一般的选择问题还允许非前缀组合。形式化目标是找到某个 $S\subseteq C$，使 $Q_{\mathrm{init}}(S)$ 与 $Q_{\mathrm{step}}(S)$ 均为 UNSAT，而不是预先要求最少的引理数。失败的初始状态或步骤状态对可为后续组合选择提供信息，但任何据此提出的新组合都必须重新接受完整的形式化检查；若后端不能可靠解释见证中的候选取值，则不能把该见证当作安全剪枝依据。

\subsection{评价目标与结果状态}

本文在固定候选池就绪后，为每个任务设置覆盖选择与证明阶段的统一墙钟预算 $B$。LLM 调用、全部形式化查询、必要的反例求值和回退搜索均消耗该预算；候选生成是各方法共用的前置阶段，不计入选择方法的 Proof Time。评价首先比较预算内完成证明的目标属性数量（Solved Properties），再比较各方法共同成功任务从候选池就绪到证明成立的墙钟时间（Proof Time）与形式化查询次数（Formal Queries）。Proof Time 包含模型推理、形式化检查及其他选择开销；Formal Queries 按实际发起的后端请求计数，包含为定位失败或求值反例而追加的查询。Solver Time、LLM Calls／Tokens 和所选引理数作为辅助指标。未证明的任务保留其终止状态与已消耗时间，并计入完成率的分母；其已耗时间不作为成功证明时间纳入上述成本均值。

结果状态须反映证明证据的强度：\texttt{PROVED} 表示某个组合的两项归纳查询均已被证明 UNSAT。若发现从初始状态可达的 $\neg p$ 反例，则目标属性被违反；归纳步骤的 SAT 见证只否决当前组合，不能单独说明目标属性被违反。

\texttt{NO\_PROOF\_IN\_POOL} 仅在固定池搜索完整且可靠、又无有效组合时返回；它不表示目标属性不安全。预算耗尽、检查未决及工具故障分别记录为 \texttt{TIMEOUT}、\texttt{UNKNOWN} 与 \texttt{ERROR}，不得与证明成功或可达反例混同。
